\section{Konstrukcja Carathéodory'ego i mierzalność w sensie Lebesgue'a}
Zewnętrzna miara Lebesgue'a posiada znakomite własności geometryczne — pokrywa się z objętością kostek i jest niezmiennicza na przesunięcia. Posiada jednak poważną wadę: z faktu, że dwa zbiory są rozłączne, nie wynika wcale, że miara ich sumy jest sumą ich miar (zachodzi jedynie podaddytywność). Mamy zatem funkcję, która mierzy wszystko, ale nie zawsze addytywnie. Powstaje fundamentalny dylemat: jak wyselekcjonować rodzinę takich podzbiorów, dla których postulowana przeliczalna addytywność zostanie w pełni odzyskana? Rozwiązanie tego problemu podał Constantin Carathéodory, wprowadzając warunek czysto geometrycznego rozcinania przestrzeni.

\subsection{Warunek mierzalności i twierdzenie Carathéodory'ego} \begin{defi}[Warunek Carathéodory'ego] Niech $\phi \colon \mathcal{P}(X) \to [0, +\infty]$ będzie miarą zewnętrzną na zbiorze $X$. Zbiór $A \subseteq X$ nazywamy \textbf{mierzalnym względem $\phi$ w sensie Carathéodory'ego} (lub krótko: \textbf{$\phi$-mierzalnym}), jeżeli dla każdego zbioru testowego $T \subseteq X$ spełniony jest warunek:
    \begin{equation}\label{eq:caratheodory_condition}
        \phi(T) = \phi(T \cap A) + \phi(T \setminus A).
    \end{equation}
    Rodzinę wszystkich zbiorów $\phi$-mierzalnych oznaczamy przez $\mathcal{M}_\phi$.
\end{defi}
\begin{uwaga}
    Dla dowolnych zbiorów $T, A \subseteq X$ zachodzi inkluzja $T \subseteq (T \cap A) \cup (T \setminus A)$, skąd z podaddytywności miary zewnętrznej zawsze mamy nierówność
    \begin{equation*}
        \phi(T) \mlr \phi(T \cap A) + \phi(T \setminus A).
    \end{equation*}
    Zatem aby wykazać, że zbiór $A$ jest $\phi$-mierzalny, wystarczy sprawdzić nierówność przeciwną:
    \begin{equation}\label{eq:carath_reduced}
        \forall_{T \subseteq X} \phi(T) < +\infty \implies \phi(T) \wlr \phi(T \cap A) + \phi(T \setminus A).
    \end{equation}
\end{uwaga}

\begin{tw}[Twierdzenie Carathéodory'ego]\label{tw:caratheodory}
    Niech $\phi$ będzie miarą zewnętrzną na zbiorze $X$. Wówczas:
    \begin{enumerate}
        \item[(i)] Rodzina $\mathcal{M}_\phi$ jest $\sigma$-algebrą podzbiorów $X$.
        \item[(ii)] Obcięcie $\mu := \phi|_{\mathcal{M}_\phi}$ jest miarą przeliczalnie addytywną i zupełną na $\mathcal{M}_\phi$.
    \end{enumerate}
\end{tw}
\begin{dow} Wykażemy twierdzenie krok po kroku.
    \paragraph{Krok 1: $\mathcal{M}_\phi$ jest algebrą zbiorów.} Dla zbioru pustego oraz dowolnego $T \subseteq X$ mamy:

    \begin{equation*} \phi(T \cap \emptyset) + \phi(T \setminus \emptyset) = \phi(\emptyset) + \phi(T) = 0 + \phi(T) = \phi(T),
    \end{equation*}
    zatem $\emptyset \in \mathcal{M}_\phi$. Warunek \eqref{eq:caratheodory_condition} jest symetryczny ze względu na dopełnienie, ponieważ $T \setminus A = T \cap A^c$ oraz $T \cap A = T \setminus A^c$.
    Stąd jeśli $A \in \mathcal{M}_\phi$, to $A^c \in \mathcal{M}_\phi$.
    Niech teraz $A, B \in \mathcal{M}_\phi$ i niech $T \subseteq X$ będzie dowolnym zbiorem testowym. Korzystając z mierzalności $A$, a następnie $B$ dla zbioru testowego $T \setminus A$, otrzymujemy:
    \begin{align}
        \phi(T) & = \phi(T \cap A) + \phi(T \setminus A)                                               \\
                & = \phi(T \cap A) + \phi((T \setminus A) \cap B) + \phi((T \setminus A) \setminus B).
    \end{align} Zauważmy, że:
    \begin{equation*}
        (T \cap A) \cup ((T \setminus A) \cap B) = T \cap (A \cup B) \quad \text{oraz} \quad (T \setminus A) \setminus B = T \setminus (A \cup B).
    \end{equation*}
    Z podaddytywności $\phi$ mamy:
    \begin{equation*}
        \phi(T \cap (A \cup B)) \mlr \phi(T \cap A) + \phi((T \setminus A) \cap B). \end{equation*}
    Wstawiając to do powyższej równości, otrzymujemy:
    \begin{equation*}
        \phi(T) \ge \phi(T \cap (A \cup B)) + \phi(T \setminus (A \cup B)),
    \end{equation*}
    co dowodzi, że $A \cup B \in \mathcal{M}_\phi$. Zatem $\mathcal{M}_\phi$ jest algebrą zbiorów.

    \paragraph{Krok 2: Skończona addytywność warunku rozcinania.}
    Niech $A, B \in \mathcal{M}_\phi$ będą zbiorami rozłącznymi ($A \cap B = \emptyset$). Podstawiając w warunku mierzalności $A$ zbiór testowy $T \cap (A \cup B)$, dostajemy:
    \begin{equation}\label{eq:disjoint_split}
        \phi(T \cap (A \cup B)) = \phi(T \cap A) + \phi(T \cap B).
    \end{equation}
    Przez prostą indukcję: jeżeli zbiory $A_1, \dots, A_m \in \mathcal{M}_\phi$ są parami rozłączne, to dla każdego $T \subseteq X$:
    \begin{equation}\label{eq:finite_disjoint_split}
        \phi\left(T \cap \bigcup_{j=1}^m A_j\right) = \sum_{j=1}^m \phi(T \cap A_j).
    \end{equation}

    \paragraph{Krok 3: Zamkniętość $\mathcal{M}_\phi$ na przeliczalne sumy.}
    Niech $\{A_j\}_{j=1}^\infty \subseteq \mathcal{M}_\phi$ będzie ciągiem parami rozłącznych zbiorów mierzalnych i niech $A = \bigcup_{j=1}^\infty A_j$. Dla dowolnego $m \in \N$, kładąc $B_m = \bigcup_{j=1}^m A_j \in \mathcal{M}_\phi$, mamy dla dowolnego $T \subseteq X$:
    \begin{align}
        \phi(T) & = \phi(T \cap B_m) + \phi(T \setminus B_m)                                                                      \\
                & = \sum_{j=1}^m \phi(T \cap A_j) + \phi(T \setminus B_m) \quad \text{(na mocy \eqref{eq:finite_disjoint_split})} \\
                & \wlr \sum_{j=1}^m \phi(T \cap A_j) + \phi(T \setminus A),
    \end{align}
    gdzie w ostatnim kroku skorzystaliśmy z faktu, że $B_m \subseteq A \implies T \setminus A \subseteq T \setminus B_m$ i monotoniczności $\phi$.
    Ponieważ nierówność ta zachodzi dla każdego $m$, przechodząc do granicy $m \to \infty$, uzyskujemy:
    \begin{equation*}
        \phi(T) \ge \sum_{j=1}^\infty \phi(T \cap A_j) + \phi(T \setminus A) \ge \phi(T \cap A) + \phi(T \setminus A),
    \end{equation*}
    gdzie ostatnia nierówność wynika z przeliczalnej podaddytywności $\phi$. To dowodzi, że $A \in \mathcal{M}_\phi$. Każdą przeliczalną sumę zbiorów $\bigcup_{j=1}^\infty E_j$ z algebry można zapisać jako sumę zbiorów parami rozłącznych: $A_1 = E_1$, $A_j = E_j \setminus \bigcup_{k=1}^{j-1} E_k \in \mathcal{M}_\phi$.
    Wobec tego $\mathcal{M}_\phi$ jest zamknięta na przeliczalne sumy, czyli jest $\sigma$-algebrą.

    \paragraph{Krok 4: $\phi|_{\mathcal{M}_\phi}$ jest miarą zupełną.}
    Kładąc w równości \eqref{eq:finite_disjoint_split} zbiór testowy $T = A = \bigcup_{j=1}^\infty A_j$, otrzymujemy:
    \begin{equation*}
        \phi\left(\bigcup_{j=1}^\infty A_j\right) \wlr \sum_{j=1}^\infty \phi(A_j).
    \end{equation*}
    Nierówność przeciwna wynika z podaddytywności miary zewnętrznej. Zatem $\phi$ jest przeliczalnie addytywna na $\mathcal{M}_\phi$. Pozostaje wykazać zupełność. Załóżmy, że $Z \subseteq X$ spełnia $\phi(Z) = 0$ oraz niech $E \subseteq Z$. Z monotoniczności $\phi$ wynika, że $\phi(E) = 0$. Wtedy dla dowolnego $T \subseteq X$:
    \begin{equation*}
        \phi(T \cap E) \mlr \phi(E) = 0 \implies \phi(T \cap E) = 0,
    \end{equation*} skąd:
    \begin{equation*}
        \phi(T \cap E) + \phi(T \setminus E) = 0 + \phi(T \setminus E) \mlr \phi(T).
    \end{equation*}
    Zatem każdy podzbiór zbioru o mierze zero jest $\phi$-mierzalny, co kończy dowód zupełności.
\end{dow}
\subsection{Mierzalność zbiorów borelowskich i definicja miary $\lambda^n$}
Aby wykazać, że rodzina zbiorów mierzalnych jest wystarczająco bogata (w szczególności zawiera wszystkie zbiory borelowskie), zbadamy związek miary zewnętrznej z metryką przestrzeni.
\begin{defi}[Metryczna miara zewnętrzna]
    Niech $(X, d)$ będzie przestrzenią metryczną. Miarę zewnętrzną $\phi \colon \mathcal{P}(X) \to [0, +\infty]$ nazywamy \textbf{metryczną miarą zewnętrzną}, jeżeli jest ona addytywna na zbiorach odseparowanych metrycznie, to znaczy:
    \begin{equation*}
        \forall A, B \subseteq X \colon \quad \dist(A, B) > 0 \implies \phi(A \cup B) = \phi(A) + \phi(B),
    \end{equation*}
    gdzie $\dist(A, B) := \inf\{d(x, y) \colon x \in A, \, y \in B\}$ (przyjmując $\dist(A, B) = +\infty$, gdy $A = \emptyset$ lub $B = \emptyset$).
\end{defi}
\begin{tw}[Mierzalność borelowska dla miar metrycznych]\label{tw:metric_borel}
    Niech $(X, d)$ będzie przestrzenią metryczną, a $\phi$ metryczną miarą zewnętrzną na $X$. Wówczas każdy zbiór borelowski jest $\phi$-mierzalny, to znaczy:
    \begin{equation*}
        \Borel{X} \subseteq \mathcal{M}_\phi.
    \end{equation*}
\end{tw}

\paragraph{Idea dowodu.}
Ponieważ $\mathcal{M}_\phi$ jest $\sigma$-algebrą, wystarczy sprawdzić warunek Carathéodory'ego dla zbiorów otwartych $G$ i zbiorów testowych $T$ o skończonej mierze. Zbiór $G$ przybliżamy od wewnątrz wstępującym ciągiem zbiorów $G_k := \{x \in G \colon \dist(x, G^c) > \frac{1}{k}\}$, których suma jest równa $G$. Ponieważ $\dist(G_k, G^c) \wlr \frac{1}{k} > 0$, z metryczności $\phi$ dostajemy $\phi(T \cap G_k) + \phi(T \setminus G) \mlr \phi(T)$. Pozostaje pokazać, że $\phi(T \cap (G \setminus G_k)) \to 0$; w tym celu dzieli się $G \setminus G_k$ na ,,warstwy'' $\{x \colon \frac{1}{j+1} < \dist(x, G^c) \mlr \frac{1}{j}\}$, $j \wlr k$, i korzysta z tego, że warstwy o numerach różniących się co najmniej o $2$ są odseparowane metrycznie, co daje zbieżność szeregu miar ich przecięć z $T$. Pełny dowód można znaleźć np.\ w~\cite{Plebanek}.

\begin{tw}[Metryczność miary $\lambda^{*n}$]

    Zewnętrzna miara Lebesgue'a $\lambda^{*n}$ w $\Rn$ jest metryczną miarą zewnętrzną.

\end{tw}
\begin{dow} Niech $A, B \subseteq \Rn$ spełniają warunek $\rho := \dist(A, B) > 0$. Z podaddytywności miary zewnętrznej mamy $\lambda^{*n}(A \cup B) \mlr \lambda^{*n}(A) + \lambda^{*n}(B)$. Należy wykazać nierówność przeciwną. Możemy bez zmniejszenia ogólności założyć, że $\lambda^{*n}(A \cup B) < +\infty$. Ustalmy dowolne $\varepsilon > 0$. Zgodnie z własnością pokryć kostkami o ograniczonej średnicy, istnieje przeliczalne pokrycie zbioru $A \cup B$ kostkami domkniętymi $\{P_j\}_{j=1}^\infty$ takimi, że:
    \begin{equation*}
        A \cup B \subseteq \bigcup_{j=1}^\infty P_j, \quad \diam(P_j) < \frac{\rho}{3} \quad \text{dla każdego } j \in \N,
    \end{equation*}
    oraz
    \begin{equation*}
        \sum_{j=1}^\infty V^n(P_j) \mlr \lambda^{*n}(A \cup B) + \varepsilon.
    \end{equation*}
    Ponieważ $\diam(P_j) < \frac{\rho}{3}$, żadna kostka $P_j$ nie może przecinać jednocześnie zbioru $A$ i zbioru $B$ (gdyby istniały $x \in P_j \cap A$ oraz $y \in P_j \cap B$, to $d(x, y) \mlr \diam(P_j) < \frac{\rho}{3} < \rho$, co przeczy definicji $\rho = \dist(A, B)$). Możemy zatem podzielić rodzinę wskaźników $\N$ na trzy parami rozłączne podzbiory:
    \begin{equation*}
        J_A = \{j \in \N \colon P_j \cap A \neq \emptyset\}, \quad J_B = \{j \in \N \colon P_j \cap B \neq \emptyset\}, \quad J_0 = \N \setminus (J_A \cup J_B).
    \end{equation*}
    Wówczas rodzina $\{P_j\}_{j \in J_A}$ stanowi pokrycie zbioru $A$, a rodzina $\{P_j\}_{j \in J_B}$ stanowi pokrycie zbioru $B$. Z definicji miary zewnętrznej jako infimum sum objętości otrzymujemy:
    \begin{equation*}
        \lambda^{*n}(A) \mlr \sum_{j \in J_A} V^n(P_j) \quad \text{oraz} \quad \lambda^{*n}(B) \mlr \sum_{j \in J_B} V^n(P_j).
    \end{equation*}
    Sumując obie nierówności:
    \begin{equation*}
        \lambda^{*n}(A) + \lambda^{*n}(B) \mlr \sum_{j \in J_A} V^n(P_j) + \sum_{j \in J_B} V^n(P_j) \mlr \sum_{j=1}^\infty V^n(P_j) \mlr \lambda^{*n}(A \cup B) + \varepsilon.
    \end{equation*}
    Z dowolności $\varepsilon > 0$ wynika, że $\lambda^{*n}(A) + \lambda^{*n}(B) \mlr \lambda^{*n}(A \cup B)$, co daje żądaną równość.
\end{dow}
\begin{wn}[Mierzalność zbiorów borelowskich] Z Twierdzenia~\ref{tw:metric_borel} i metryczności $\lambda^{*n}$ wynika inkluzja:
    \begin{equation*}
        \Borel{\Rn} \subseteq \mathcal{M}_{\lambda^{*n}}.
    \end{equation*}
    W szczególności każda kostka, kula oraz każdy zbiór otwarty i domknięty w $\Rn$ jest mierzalny w sensie Carathéodory'ego względem zewnętrznej miary Lebesgue'a.
\end{wn}
\begin{defi}[Miara Lebesgue'a i $\sigma$-algebra $\Leb{\Rn}$]
    $\sigma$-algebrę $\mathcal{M}_{\lambda^{*n}}$ zbiorów spełniających warunek Carathéodory'ego względem zewnętrznej miary Lebesgue'a oznaczamy symbolem $\mathcal{L}(\Rn)$ (lub $\mathcal{L}^n$) i nazywamy \textbf{$\sigma$-algebrą zbiorów mierzalnych w sensie Lebesgue'a}. Obcięcie miary zewnętrznej do tej $\sigma$-algebry:
    \begin{equation*}
        \lambda^n := \lambda^{*n}|_{\Leb{\Rn}}
    \end{equation*}
    nazywamy \textbf{$n$-wymiarową miarą Lebesgue'a} w $\Rn$. Trójka $(\Rn, \Leb{\Rn}, \lambda^n)$ stanowi zupełną przestrzeń z miarą.
\end{defi}
